\documentclass[11pt]{article}

\usepackage[T1]{fontenc}
\usepackage{lmodern}
\usepackage{geometry}
\geometry{margin=1in}

\usepackage{amsmath,amssymb,amsthm,mathtools}
\usepackage{enumitem}
\usepackage{hyperref}

% ---------- Macros ----------
\newcommand{\ZZ}{\mathbb{Z}}
\newcommand{\QQ}{\mathbb{Q}}
\newcommand{\Sym}{\mathrm{Sym}}
\newcommand{\NF}{\mathrm{NF}}
\newcommand{\GL}{\mathrm{GL}}
\newcommand{\rank}{\mathrm{rank}}
\newcommand{\sign}{\mathrm{sign}}
\newcommand{\lex}{\mathrm{lex}}
\newcommand{\proj}{\pi}
\newcommand{\bracks}[1]{\llbracket #1 \rrbracket}

\newcommand{\cfg}[2]{\langle #1,#2\rangle}

\newcommand{\Rule}[1]{\textsc{(#1)}}

% ---------- Title ----------
\title{Operator-Word Calculus on \texorpdfstring{$(Q,p)$}{(Q,p)}:\\
Slides, Blow-Ups, Hyperbolic Stabilization, and Rewrite-to-\texorpdfstring{$\NF(Q,p)$}{NF(Q,p)}}
\author{Ryan Van Gelder}
\date{2025}

\begin{document}
\maketitle

\begin{abstract}
We formalize a state-based rewriting logic for an operator-word calculus acting on configurations
$(Q,p)$, where $Q$ is an integral symmetric form and $p$ is an integral refinement vector.
The generator set consists of slides $S_{ij}^{\pm}$ (elementary unimodular basis changes),
blow up/down moves $B_{\pm}, B_{\pm}^{-1}$, and hyperbolic stabilization/cancellation $H,H^{-1}$.
A deterministic rewrite strategy produces a canonical normal form $\NF(Q,p)$ that is invariant
under the generated equivalences. We also indicate a specialization to 8D plumbings bounding Milnor
7-spheres, making admissibility tests and the canonical $p$-data explicit.
\end{abstract}

\paragraph{Background pointers.}
Exotic $7$-spheres and $\Theta_7$: \cite{Mil56,KM63}. Smoothing theory: \cite{KS77}. Kirby calculus:
\cite{Kir78,GS99}.

% ============================================================
\section{Configurations and semantics}

\subsection{Configuration space}
A configuration is a pair
\[
(Q,p),\qquad Q\in \Sym_n(\ZZ),\quad p\in \ZZ^n.
\]
Interpretation: $Q$ is a lattice pairing (e.g.\ an intersection form in a chosen basis) and $p$ is
extra structure (e.g.\ a characteristic/Pontryagin-refinement datum in that basis).

\subsection{Congruence (``gauge'') action}
A unimodular basis change $U\in \GL_n(\ZZ)$ acts by
\[
(Q,p)\longmapsto (U^T Q U,\; U^{-1}p).
\]
Slides are elementary generators for this action (handle slides in Kirby calculus; see \cite{Kir78,GS99}).

% ============================================================
\section{Generator operators (raw moves)}

\subsection{Slides \texorpdfstring{$S_{ij}^{\varepsilon}$}{Sij}}
For $i\neq j$ and $\varepsilon\in\{+1,-1\}$, let
\[
U_{ij}(\varepsilon)=I+\varepsilon E_{ij}\in \GL_n(\ZZ),
\]
where $E_{ij}$ has a single $1$ in position $(i,j)$.

\paragraph{Definition (Slide).}
\[
S_{ij}^{\varepsilon}(Q,p)=\big(U_{ij}(\varepsilon)^T Q\, U_{ij}(\varepsilon),\; U_{ij}(\varepsilon)^{-1}p\big).
\]
Slides are total and invertible: $(S_{ij}^{\varepsilon})^{-1}=S_{ij}^{-\varepsilon}$.

\subsection{Blow-up/down \texorpdfstring{$B_\sigma$}{B}}
For $\sigma\in\{+1,-1\}$ define the blow-up
\[
B_\sigma(Q,p)=\big(Q\oplus[\sigma],\; p\oplus 0\big).
\]
This is total.

\paragraph{Blow-down $B_\sigma^{-1}$ (partial).}
It is admissible only if $(Q,p)$ contains an orthogonal primitive $\sigma$-vector (precisely in \S\ref{sec:blowdown}).

\subsection{Hyperbolic stabilization/cancellation \texorpdfstring{$H$}{H}}
Let $H_2=\begin{pmatrix}0&1\\1&0\end{pmatrix}$. Define
\[
H(Q,p)=\big(Q\oplus H_2,\; p\oplus(0,0)\big).
\]
Cancellation $H^{-1}$ is partial: it is admissible only if $(Q,p)$ contains an orthogonal hyperbolic plane (precisely in \S\ref{sec:hypsplit}).

% ============================================================
\section{Rewriting logic: explicit rules on configurations}

We write rewrite rules as
\[
\cfg{Q}{p}\;\Rightarrow\;\cfg{Q'}{p'}
\]
possibly under an admissibility predicate.

\subsection{Slide rule (always admissible)}
\[
\Rule{SLIDE}\qquad
\frac{}{\cfg{Q}{p}\Rightarrow \cfg{U^T Q U}{U^{-1}p}}
\]
where $U=U_{ij}(\varepsilon)=I+\varepsilon E_{ij}$.

\subsection{Blow-up rule (always admissible)}
\[
\Rule{BLOWUP}\qquad
\frac{}{\cfg{Q}{p}\Rightarrow \cfg{Q\oplus[\sigma]}{p\oplus 0}}
\qquad \sigma\in\{\pm1\}.
\]

\subsection{Blow-down rule (partial)}\label{sec:blowdown}
We formalize blow-down by finding a unimodular $U$ such that
\[
U^T Q U = Q'\oplus[\sigma].
\]
Let $\proj$ be projection onto the first $n-1$ coordinates.

\paragraph{Predicate $\mathrm{Split}_\sigma(Q)$.}
There exists primitive $v\in \ZZ^n$ and $U\in \GL_n(\ZZ)$ with $Ue_n=v$ and $U^T Q U=Q'\oplus[\sigma]$.

\[
\Rule{BLOWDOWN}\qquad
\frac{\mathrm{Split}_\sigma(Q)\ \wedge\ U^TQU=Q'\oplus[\sigma]}{\cfg{Q}{p}\Rightarrow \cfg{Q'}{\proj(U^{-1}p)}}
\]
In rewriting-to-normal-form we do not allow arbitrary $U$; we choose a canonical witness (see \S\ref{sec:canonwitness}).

\subsection{Hyperbolic add rule (always admissible)}
\[
\Rule{HYPADD}\qquad
\frac{}{\cfg{Q}{p}\Rightarrow \cfg{Q\oplus H_2}{p\oplus(0,0)}}.
\]

\subsection{Hyperbolic cancel rule (partial)}\label{sec:hypsplit}
\paragraph{Predicate $\mathrm{HypSplit}(Q)$.}
There exist primitive $u,v\in\ZZ^n$ and $U\in \GL_n(\ZZ)$ so that $Ue_{n-1}=u$, $Ue_n=v$, and
\[
U^T Q U = Q'\oplus H_2.
\]
Equivalently, in the original basis: $u^TQu=0$, $v^TQv=0$, $u^TQv=1$, and the span splits off orthogonally.

Let $\proj$ project away the last two coordinates. Then
\[
\Rule{HYPCANCEL}\qquad
\frac{\mathrm{HypSplit}(Q)\ \wedge\ U^TQU=Q'\oplus H_2}{\cfg{Q}{p}\Rightarrow \cfg{Q'}{\proj(U^{-1}p)}}.
\]
Again, normalization uses a canonical choice of $(u,v,U)$ (see \S\ref{sec:canonwitness}).

% ============================================================
\section{Deterministic rewrite-to-normal-form strategy}

The central idea is: \emph{words are not normalized syntactically; states are normalized deterministically.}

Define a deterministic normalization map
\[
\mathsf{N}:\ (Q,p)\longmapsto \NF(Q,p).
\]
We obtain a normalized step relation
\[
(Q,p)\xRightarrow{g}\mathsf{N}\big(g(Q,p)\big)
\]
for any generator $g$ admissible at $(Q,p)$.

\subsection{Canonical witness selection for partial reductions}\label{sec:canonwitness}
To make $\mathsf N$ deterministic, every time a partial rule (BLOWDOWN/HYPCANCEL) is applicable, we choose the witness by a fixed minimization procedure.

\subsubsection*{Canonical blow-down choice}
Among all primitive $v$ that split off a $[\sigma]$, select the one minimizing
\[
(\|v\|_2,\ \lex(v)).
\]
Then among unimodular $U$ with $Ue_n=v$ and $U^TQU=Q'\oplus[\sigma]$, choose $U$ minimizing a deterministic tie-break key (e.g.\ lexicographic order of the last row, then of the full matrix entries).

\subsubsection*{Canonical hyperbolic cancellation choice}
Among all primitive hyperbolic pairs $(u,v)$ that split off $H_2$, select the one minimizing
\[
(\max(\|u\|_2,\|v\|_2),\ \lex(u),\ \lex(v)).
\]
Then choose a canonical completing $U$ as above.

\paragraph{Stage N1 --- Stabilization reduction.}
\begin{enumerate}[label=\arabic*.]
\item While a canonical blow-down is admissible, apply \Rule{BLOWDOWN} using the canonical witness.
\item While a canonical hyperbolic cancellation is admissible, apply \Rule{HYPCANCEL} using the canonical witness.
\item Repeat until neither applies.
\end{enumerate}
This strictly decreases rank whenever it fires, hence terminates.

\paragraph{Stage N2 --- Global congruence normalization (real-place basis fix).}
Apply a deterministic lattice reduction procedure to fix the remaining congruence freedom, e.g.:
\begin{itemize}
\item compute an LLL-reduced basis for the lattice with Gram matrix $Q$ \cite{LLL82},
\item among all LLL-valid reduced bases, choose the one that lexicographically minimizes the coordinates of $p$,
\item fix a sign convention (first nonzero entry of $p$ positive; then lex-min for remaining symmetries).
\end{itemize}
This yields a reproducible representative in the congruence class.

\paragraph{Stage N3 --- Prime-local block normalization (optional).}
(Local classification/Jordan decompositions and invariants of integral quadratic forms: see \cite{OM63,Ser73,CS99}.)
For each selected prime $\ell$ (and optionally depth $r$):
\begin{enumerate}[label=\arabic*.]
\item Compute the $\ell$-adic Jordan decomposition of $(Q,p)\otimes \ZZ_\ell$:
\[
(Q,p)\otimes\ZZ_\ell \cong \bigoplus_{e\ge 0}\ell^e\,(Q_{\ell,e},\,p_{\ell,e}).
\]
\item If using a projector $\Pi_{\ell^r}$, truncate to $0\le e<r$.
\item Sort blocks by a canonical key, e.g.
\[
(e,\ \rank,\ \det\bmod(\ZZ_\ell^\times)^2,\ \text{Hasse/2-adic refinements},\ \text{coordinates of }p_{\ell,e}).
\]
\end{enumerate}

The result is $\NF(Q,p)$. If only a stable normal form is needed, one may stop at Stage N2.

% ============================================================
\section{Operator words, evaluation, and equivalence}

\subsection{Words}
A word is a sequence of generator symbols
\[
w=g_m\cdots g_2 g_1
\]
subject to admissibility of partial inverses at runtime.

\subsection{Normalized evaluation}
Define
\[
\mathrm{EvalNF}(w;Q,p) := \mathsf{N}\big(\bracks{w}(Q,p)\big).
\]
Two words are equivalent on input $(Q,p)$ if they yield the same normal form.

\subsection{``Resonance = 1'' criterion (move invariance)}
Replace any scoring notion by strict equality of canonical outputs:
\[
\mathrm{EvalNF}(w_1;Q,p) = \mathrm{EvalNF}(w_2;Q,p).
\]

% ============================================================
\section{Specialization: 8D plumbing coboundaries of Milnor 7-spheres}\label{sec:plumbing}

This section makes $(Q,p)$ explicitly computable from an 8D plumbing of disk bundles over $S^4$ whose boundary is a homotopy 7-sphere (Milnor-type).

\subsection{Plumbing input}
A plumbing description consists of:
\begin{itemize}
\item a finite graph $\Gamma=(V,E)$;
\item for each vertex $v\in V$, an oriented $D^4$-bundle over $S^4$ with structure group $SO(4)$, specified by bundle parameters $(a_v,b_v)\in \ZZ\oplus \ZZ$ (a standard $\pi_3(SO(4))\cong \ZZ\oplus\ZZ$ coordinate choice; see \cite{DW59,CE03});
\item each edge represents a plumbing along fiber/base disks with a chosen sign convention.
\end{itemize}

Let $n=|V|$ and choose the basis $\{x_v\}_{v\in V}\subset H^4(W;\ZZ)$ given by the core $S^4$ of each bundle.

\subsection{Constructing the intersection matrix $Q$}
Define $Q\in \Sym_n(\ZZ)$ by:
\begin{itemize}
\item Diagonal:
\[
Q_{vv}=e_v,\qquad e_v:=a_v+b_v \quad \text{(Euler number in this convention).}
\]
\item Off-diagonal: for $v\neq w$,
\[
Q_{vw} = \sum_{\text{edges }(v,w)} \epsilon_{vw},
\]
where $\epsilon_{vw}\in\{\pm 1\}$ is the plumbing sign (often all $+1$ in standard Milnor plumbings).
\end{itemize}
This is the usual ``diagonal = self-intersection, edges = intersection'' plumbing matrix pattern, in degree $4$ for an 8-manifold.

\subsection{Constructing the refinement vector $p$}
In the same basis $\{x_v\}$ define $p\in \ZZ^n$ by
\[
(p)_v := a_v-b_v
\quad\text{(a standard choice compatible with $p_1$ conventions up to overall factors).}
\]
\paragraph{Spin constraint / admissibility.}
To be in the spin setting, require the appropriate mod $2$ parity condition on the vertex data (equivalently $w_2(W)=0$); in many plumbing conventions this forces each Euler number $e_v$ to be even.

\subsection{A boundary scalar from $(Q,p)$ (Eells--Kuiper type)}
Once $(Q,p)$ is formed for the coboundary $W^8$, define:
\begin{itemize}
\item Signature $\sigma(W)=\sign(Q)$;
\item Quadratic value (over $\QQ$):
\[
p_W^2 := p^T Q^{-1}p\in \QQ.
\]
\end{itemize}
Then a boundary invariant (for homotopy 7-spheres in the relevant class) is built from
\[
\mu \equiv \frac{p_W^2-\sigma(W)}{224}\pmod{1},
\]
up to convention choices (Eells--Kuiper invariant \cite{EK62}; modern classification context \cite{CN19}).
This makes the $2$- and $7$-primary structure explicit via the denominator.

\subsection{Admissibility tests in plumbing cases (explicit)}
In plumbing-derived $Q$, admissibility for cancellations often becomes concrete.

\paragraph{Blow-down test ($\pm 1$-sphere split).}
Search for primitive $v$ with
\[
v^TQv=\pm 1.
\]
If found, attempt to complete $v$ to a unimodular basis and check if the last basis vector is orthogonal to the others (i.e.\ off-diagonal last row/column can be eliminated by slides).

\paragraph{Hyperbolic split test.}
Search for primitive $u,v$ with
\[
u^TQu=0,\quad v^TQv=0,\quad u^TQv=1,
\]
and verify orthogonal split by completing $(u,v)$ to a unimodular basis and reducing cross-terms via slides.

These tests are exactly what Stage N1 requires.

% ============================================================
\section{Compact rewriting-logic summary}

\subsection{Rule schema list}
\begin{itemize}
\item \Rule{SLIDE} $\cfg{Q}{p}\Rightarrow \cfg{U^TQU}{U^{-1}p}$ for $U=I\pm E_{ij}$.
\item \Rule{BLOWUP} $\cfg{Q}{p}\Rightarrow \cfg{Q\oplus[\sigma]}{p\oplus 0}$.
\item \Rule{BLOWDOWN} $\cfg{Q}{p}\Rightarrow \cfg{Q'}{\proj(U^{-1}p)}$ if $U^TQU=Q'\oplus[\sigma]$ (canonical witness).
\item \Rule{HYPADD} $\cfg{Q}{p}\Rightarrow \cfg{Q\oplus H_2}{p\oplus(0,0)}$.
\item \Rule{HYPCANCEL} $\cfg{Q}{p}\Rightarrow \cfg{Q'}{\proj(U^{-1}p)}$ if $U^TQU=Q'\oplus H_2$ (canonical witness).
\end{itemize}

\subsection{Deterministic strategy}
Define $\NF(Q,p)=\mathsf N(Q,p)$ by:
\begin{enumerate}[label=\arabic*.]
\item Apply canonical \Rule{BLOWDOWN}/\Rule{HYPCANCEL} until stuck.
\item Apply deterministic congruence normalization (LLL + lex tie-break on $p$).
\item Optionally apply prime-local block normalization (Jordan decomposition + sorted blocks), including truncations $\Pi_{\ell^r}$ when desired.
\end{enumerate}

% ============================================================
\section{Suggested minimum viable experiment (Milnor 7-spheres)}
\paragraph{Input.} A plumbing graph $\Gamma$ with vertex parameters $(a_v,b_v)$ and edge signs.

\paragraph{Compute.}
\begin{enumerate}[label=\arabic*.]
\item $(Q,p)$ via \S\ref{sec:plumbing}.
\item $\NF(Q,p)$ via \S4.
\item Prime-tier fingerprints via Stage N3 (especially $\ell=2$ and $\ell=7$).
\item Scalar $\mu$ and its $2$- and $7$-primary projections.
\end{enumerate}

\paragraph{Validate invariance.}
Randomly apply sequences of slides (and optional stabilizations followed by normalization) and confirm $\NF(Q,p)$ is unchanged.

% ============================================================
\section{Notes on conventions}
The precise mapping from vertex parameters $(a_v,b_v)$ to $(e_v,(p)_v)$ depends on orientation/normalization conventions for $SO(4)$-bundles over $S^4$. The rewrite calculus is invariant under any consistent choice, as long as:
\begin{enumerate}[label=\arabic*.]
\item $Q$ is integral symmetric and transforms by congruence under slides,
\item $p$ transforms as $p\mapsto U^{-1}p$, and
\item stabilizations append zeros to $p$.
\end{enumerate}
If you adopt a different convention, record it once and keep it fixed; $\NF(Q,p)$ remains well-defined for that convention.

% ============================================================
\begin{thebibliography}{99}

\bibitem{MOC25}
\textsc{Multiplicity Operator Calculus (MOC)} (internal note/manuscript), Nov.\ 5, 2025 (unpublished).

\bibitem{PMDM25}
\textsc{PMDM: Exotic Spheres direction} (internal note/manuscript), 2025 (unpublished).

\bibitem{CN19}
D.~Crowley and J.~Nordstr{\"o}m,
\newblock The classification of 2-connected 7-manifolds,
\newblock \emph{Proc.\ London Math.\ Soc.}\ (3) \textbf{119} (2019), no.~1, 1--54.

\bibitem{CE03}
D.~Crowley and C.~M.~Escher,
\newblock A classification of $S^3$-bundles over $S^4$,
\newblock \emph{Differential Geom.\ Appl.}\ \textbf{18} (2003), no.~3, 363--380.

\bibitem{DW59}
A.~Dold and H.~Whitney,
\newblock Classification of oriented sphere bundles over a 4-complex,
\newblock \emph{Ann.\ of Math.}\ (2) \textbf{69} (1959), no.~3, 667--677.

\bibitem{EK62}
J.~Eells, Jr.\ and N.~H.~Kuiper,
\newblock An invariant for certain smooth manifolds,
\newblock \emph{Annali di Matematica Pura ed Applicata}\ (4) \textbf{60} (1962), 93--110.

\bibitem{GS99}
R.~E.~Gompf and A.~I.~Stipsicz,
\newblock \emph{4-Manifolds and Kirby Calculus},
\newblock Graduate Studies in Mathematics, Vol.\ 20, AMS, 1999.

\bibitem{Kir78}
R.~Kirby,
\newblock A calculus for framed links in $S^3$,
\newblock \emph{Invent.\ Math.}\ \textbf{45} (1978), 35--56.

\bibitem{KM63}
M.~A.~Kervaire and J.~W.~Milnor,
\newblock Groups of homotopy spheres. I,
\newblock \emph{Ann.\ of Math.}\ (2) \textbf{77} (1963), no.~3, 504--537.

\bibitem{LLL82}
A.~K.~Lenstra, H.~W.~Lenstra Jr., and L.~Lov{\'a}sz,
\newblock Factoring polynomials with rational coefficients,
\newblock \emph{Math.\ Ann.}\ \textbf{261} (1982), 515--534.

\bibitem{Mil56}
J.~Milnor,
\newblock On manifolds homeomorphic to the 7-sphere,
\newblock \emph{Ann.\ of Math.}\ (2) \textbf{64} (1956), no.~2, 399--405.

\bibitem{OM63}
O.~T.~O'Meara,
\newblock \emph{Introduction to Quadratic Forms},
\newblock Springer, 1963.

\bibitem{Ser73}
J.-P.~Serre,
\newblock \emph{A Course in Arithmetic},
\newblock Springer, 1973.

\bibitem{CS99}
J.~H.~Conway and N.~J.~A.~Sloane,
\newblock \emph{Sphere Packings, Lattices and Groups}, 3rd ed.,
\newblock Springer, 1999.

\bibitem{KS77}
R.~C.~Kirby and L.~C.~Siebenmann,
\newblock \emph{Foundational Essays on Topological Manifolds, Smoothings, and Triangulations},
\newblock Princeton Univ.\ Press, 1977.

\bibitem{Wal62}
C.~T.~C.~Wall,
\newblock Classification of $(n-1)$-connected $2n$-manifolds,
\newblock \emph{Ann.\ of Math.}\ (2) \textbf{75} (1962), 163--189.

\end{thebibliography}

\vspace{1em}
\noindent\textcopyright\ 2025 Citizen Gardens. CC-NC-ND 4.0.

\end{document}
